\section{符号说明}

主要符号按后文出现顺序排列。未列入下表的中间量在首次出现处定义。坐标单位一律为米，角度单位为度，时间单位为秒。符号总数控制在正文常用对象，不把实现内部的临时变量写入本表。

\begin{center}
\begin{longtable}{p{2.2cm}p{10.8cm}}
\caption{主要符号}\\
\hline
符号 & 含义\\
\hline
\endfirsthead
\hline
符号 & 含义\\
\hline
\endhead
\hline
\endfoot
$\Omega$ & 目标圆域，$x^2+y^2\le 1800^2$，只约束源位置\\
$N$ & 干扰源总数，在线未知，$N\in\{10,\ldots,16\}$\\
$K$ & 本局已成功清除个数，$0\le K\le N$\\
$c,k$ & 测向机当前频道与本次指令频道，取值 $1$ 至 $20$\\
$p,q$ & 上一动作点与本次目标点，允许出圆\\
$G_i,g$ & 第 $i$ 个源及其位置，$g\in\Omega$\\
$r_i$ & 源 $i$ 的有效接收半径，$r_i\in[1000,1500]$\\
$\hat\theta$ & 示向度，真方位加 $[-1^\circ,1^\circ]$ 误差，两位小数\\
$\varepsilon$ & 示向度半宽，$1^\circ$；数值包络取 $1.005^\circ$\\
$\psi_i,d$ & 定向源朝向及其单位向量\\
$R$ & 前向闭楔形交集给出的定位区域\\
$\mathrm{diam}(R)$ & 定位区域直径，顶点最远对\\
$\rho$ & 最小包围圆半径（Jung 半径）\\
$J(q)$ & 第二测点排序指标，非清除证书\\
$T$ & 虚拟任务时间，含空扫与失败动作\\
$T/K$ & 平均定位清除时间；$K=0$ 时未定义\\
$P_3$ & 问题3七点发现集：原点与半径 $1200\,\mathrm{m}$ 正六边形\\
$P_4$ & 问题4正方形覆盖 SQUARE81，步长 $600\,\mathrm{m}$，共 $81$ 点\\
HEX37 & 三角晶格覆盖，步长 $800\,\mathrm{m}$，六边形半径 $3$，共 $37$ 点\\
$\delta$ & HEX37 覆盖半径，$800/\sqrt{3}$\\
$\rho_{\mathrm{th}}$ & 理论扇半径，$r_{\max}/(2\cos\varepsilon)$\\
$R_{\mathrm{core}}$ & 问题2核心可观测半径，$\max\{r_{\min}-\rho_1,1\}$\\
$y$ & 覆盖证明辅助点 $g+500d$，不进入在线策略\\
\end{longtable}
\end{center}

$T$ 与程序运行墙钟、测试窗口是三套时钟，表1第四列只报告墙钟。$K=0$ 时 $T/K$ 未定义，单元格应留空而不是填零。$P_3$ 与 $P_4$、HEX37 不共用覆盖命题：前者服务全向发现，后二者必须同时进入最小接收半径与 $180^\circ$ 开半平面。$J(q)$ 与 $\rho$ 也不共用含义：前者是外包络上的排序估计，后者是真实凸集的最小包围半径；不得把 $J\le 20$ 写成 $\rho\le 20$。

辅助点 $y$ 只出现在问题4的存在性证明，在线遍历的是有限格点。
理论扇半径 $\rho_{\mathrm{th}}$ 标定问题2的核心尺度，不等于演示外包半径 $\rho_1$，因为圆域切线外包会把理想扇略微放大：演示给出 $\rho_1\approx 750.2285094687356\,\mathrm{m}$，而 $\rho_{\mathrm{th}}\approx 750.1142460329307\,\mathrm{m}$。
覆盖半径 $\delta=800/\sqrt{3}\approx 461.8802153517\,\mathrm{m}$ 只服务 HEX37，不得拿去解释 SQUARE81 的 $300\sqrt{2}$。
符号表不列入实现内部的插入阈值、源预算参数与数值包络，那些量在首次出现处定义，并且不得冒充题面常数。
清除插入上限 $600\,\mathrm{m}$、测向插入上限 $250\,\mathrm{m}$、激进试清半径 $80\,\mathrm{m}$ 属于寻路层冻结参数，光学方格边长 $28\,\mathrm{m}$ 属于清除层几何，二者也不写入本表。
